% Source: https://artofproblemsolving.com/community/c3h1648624_set_contextual_problems
% Stable identity: aops-topic-1648624-q01
Write $W$, $H$, $P$ for the sets of girls who like women's magazines, health magazines and newspapers, so $|W|=10$, $|H|=7$, $|P|=5$ and $|W\cup H\cup P|=18$. Let $k$ be the number of girls who like only women's magazines.

\textbf{(a)} Since $10+7+5=22>18$, some girls are counted more than once, and the overlaps are not given. The two examples in (b) have different values of $k$ and satisfy every condition, so $k$ is not determined.

\textbf{(b)} \emph{Lower bound.} The $18-10=8$ girls outside $W$ all like health magazines or newspapers, and so do the $10-k$ girls in $W$ who like something besides women's magazines. These $8+(10-k)=18-k$ girls all lie in $H\cup P$, which has at most $|H|+|P|=12$ members. So $18-k\le12$, that is $k\ge6$.

\emph{Upper bound.} Clearly $k\le|W|=10$.

\emph{Both bounds occur.}
\begin{itemize}
\item $k=6$: let $4$ of the $W$-girls also like health magazines, and let the $8$ girls outside $W$ consist of $3$ who like only health magazines and $5$ who like only newspapers. Then $|H|=4+3=7$, $|P|=5$, and $6$ girls like only women's magazines.
\item $k=10$: let the $10$ $W$-girls like nothing else, and let the other $8$ girls be $4$ who like both health magazines and newspapers, $3$ who like only health magazines and $1$ who likes only newspapers. Then $|H|=4+3=7$ and $|P|=4+1=5$.
\end{itemize}
So the number of girls who like only women's magazines is at least $\boxed{6}$ and at most $\boxed{10}$, and both extremes are possible.
